% ============================================================
% TITLE BLOCK (spans both columns), ABSTRACT & INDEX TERMS
% ============================================================
% \tbd{...} marks values still to be filled in (prints red).
\providecommand{\tbd}[1]{\textcolor{red}{[#1]}}

\twocolumn[{%
    \centering
    {\huge Path Planning and Onboard Obstacle Avoidance for an\\[2pt]
    Autonomous Ground Robot among Unmapped Obstacles\par}
    \vspace{1.2em}
    {\large Chabod Masere, Jonah Menzies and Balin Carkeet\par}
    \vspace{0.4em}
    {\normalsize\itshape Group 50, School of Electrical Engineering and Robotics,
    Queensland University of Technology\par}
    \vspace{0.8em}
    {\footnotesize
    \begin{tabular}{@{}c c >{\centering\arraybackslash}m{0.5\textwidth} c@{}}
        \toprule
        \textbf{Name} & \textbf{Student no.} & \textbf{Contribution} & \textbf{Share} \\
        \midrule
        Chabod Masere & n10695141 & Collision avoidance design: re-planning and reactive layers, sensing & 33\% \\
        Jonah Menzies & n11395168 & Path planning: map inflation, A* routing, PSO smoothing, waypoint ordering & 33\% \\
        Balin Carkeet & n11336692 & Guidance: path following, waypoint capture and control & 33\% \\
        \bottomrule
    \end{tabular}\par}
    \vspace{1.6em}
}]

% Unnumbered first-page note, as on an IEEE first page
{\renewcommand{\thefootnote}{}\footnotetext{EGH446 Autonomous Systems, Major
Project Part~B. Submitted to the unit coordinator, Dr Luis Mejias, October 2026.}}

{\small\bfseries\emergencystretch=1em
\noindent\hspace{1em}\textit{Abstract}---This report presents the design and
evaluation of a navigation system that enables an autonomous ground robot to
visit a set of randomly generated waypoints in an arena containing obstacles
that are unknown in advance. Five waypoints are ordered using ant colony
optimisation over true driving distances, joined by an A* route on a map
inflated for clearance, and smoothed by particle swarm optimisation into a
path the vehicle can track. Twelve obstacles are detected only through onboard
sensing. When a detected obstacle lies near the remaining path, that part of
the path is re-planned around it, and a reactive potential field acts as a
close-range backup. Both layers are designed around a required clearance of
$0.40$~m, the vehicle radius plus a safety margin.
Across $25$ randomised missions the vehicle captured every waypoint,
completing a $108$~m mission in $117$~s on average. Measured from the
simulator's noise-free pose, it kept the required clearance in every
mission, with a margin of at least $0.20$~m: the closest passes were
$0.61$~m from an obstacle and $0.63$~m from a wall, and each waypoint
was passed within $0.24$~m. The estimator's
noise model was calibrated against the measured sensor noise, a
standard deviation of $0.10$~m and $2^{\circ}$, and the clearance
margins are sized against its measured error, $99$\% of which is below
$0.18$~m. Re-planning was essential: with it switched off, five of eight
test missions breached the clearance. A $1.5$~m/s speed limit was $26$\%
faster but left only $0.09$~m of margin, so $1.0$~m/s was kept.

\hspace{1em}\textit{Index Terms}---A* search, ant colony optimisation,
autonomous ground robots, configuration space, Kalman filter, noise
characterisation, obstacle avoidance, online re-planning, particle
swarm optimisation, path planning, potential field.
\par}
\vspace{0.6em}

% ============================================================
% INTRODUCTION
% ============================================================

\section{Introduction}
\label{sec:intro}

\lettrine{A}{utonomous} ground robots are increasingly deployed in
logistics, inspection, and industrial settings where they must travel
between goal locations without human intervention. To do this reliably,
a robot must know where it is, decide where to go next, and reach that
location without colliding with its surroundings.

Part~A of this project addressed the first two of these problems in an
obstacle-free field \cite{partA_est, partA_group}: the vehicle estimated
its state and drove straight between ordered waypoints. The subsystems
carried over from Part~A are summarised in Section~\ref{sec:guidance},
with the changes made for Part~B. Part~B adds
walls, which are known from a map, and obstacles, which are not known in
advance and can only be detected by onboard sensors as the vehicle
moves. This turns a tracking problem into a planning problem. The
waypoints must be ordered to keep the route short, a path must be found
around the walls, and that path must be revised whenever a sensor
reveals an obstacle on it.

\subsection{Task and System Overview}
\label{sec:overview}

The arena is a $52 \times 41$~m walled map containing twelve static
obstacles. The simulator treats an obstacle as a point; as the unit FAQ
(item 12) allows, each is taken to be a disc of radius $0.15$~m. The vehicle is an Ackermann-steered ground
robot of radius $0.25$~m that starts at $(2, 2)$~m. Each mission places
five waypoints at random in free space, at least $1.0$~m from every
obstacle's centre. The vehicle must visit all of them and stop on the last one,
keeping a margin of $0.15$~m between its body and every wall and
obstacle. Its centre must therefore stay at least
\begin{equation}
    d_{\text{clear}} = 0.25 + 0.15 = 0.40\ \text{m}
    \label{eq:clearance}
\end{equation}
from every wall face and obstacle surface. Throughout this report,
\emph{clearance} means this distance, from the vehicle's centre to a
wall face or an obstacle's surface, and the requirement is always
$0.40$~m. Where a distance is measured from an obstacle's centre
instead, it is stated as such: the requirement is then
$0.40 + 0.15 = 0.55$~m.

Figure~\ref{fig:architecture} shows the architecture and
Table~\ref{tab:io} the inputs, outputs and rate of each subsystem. During the mission
the state estimator supplies the state estimate
$\hat{\mathbf{x}} = [\hat{x}\ \hat{y}\ \hat{\theta}]^{\top}$, guidance
turns the path into a heading command, the reactive avoidance layer
adjusts it, and the controller converts it into speed and turn-rate
commands. The re-planner watches the object detector and replaces the
path when an obstacle is found near it.

Beyond the planning and avoidance methods taught in the unit, the
design adds four elements. (i) Waypoints are ordered by true driving
distance, the length of the A* route between each pair, not by
straight-line distance. (ii) Each A* leg is smoothed by PSO into a
B-spline that respects the vehicle's turning radius, with a fallback
that guarantees smoothing never reduces clearance. (iii) Obstacles are
remembered, filtered and confirmed before they are trusted, and the
re-planner tries a wider keep-out first, so that it routes around the
open side of an obstacle rather than through a gap beside a wall.
(iv) The estimator's noise model is measured from the simulator's
noise-free pose rather than assumed, and the planning and avoidance
margins are sized against the resulting estimate error. Each element is
compared with a simpler alternative in the sections that follow.

\begin{figure*}[t]
    \centering
    \begin{tikzpicture}[
        font=\small,
        node distance=0.9cm and 0.75cm,
        blk/.style={draw, rounded corners=2pt, minimum height=0.95cm,
                    minimum width=2.2cm, align=center, fill=black!4},
        new/.style={blk, fill=CoverGold!22},
        sen/.style={blk, minimum width=1.9cm, fill=white},
        arr/.style={-{Stealth[length=2mm]}, semithick},
        lab/.style={font=\scriptsize, inner sep=1.5pt}]

        % forward path
        \node[new] (plan) {Mission planner\\[-1pt]\scriptsize ACO, A*, PSO};
        \node[blk, right=of plan]  (guid) {Guidance};
        \node[new, right=of guid]  (avoid) {Reactive\\avoidance};
        \node[blk, right=of avoid] (ctrl) {Controller};
        \node[blk, right=of ctrl]  (veh)  {Vehicle};

        % feedback and sensing
        \node[blk, below=1.5cm of veh]   (est)  {State\\estimator};
        \node[new, below=1.5cm of guid]  (rep)  {Re-planner};
        \node[sen, below=1.5cm of avoid] (det)  {Object\\detector};
        \node[sen, above=0.8cm of avoid] (lid)  {Lidar};

        \draw[arr] (plan) -- node[lab, above] {path} (guid);
        \draw[arr] (guid) -- node[lab, above] {$\theta_{\text{cmd}}$} (avoid);
        \draw[arr] (avoid) -- node[lab, above] {$\theta_{\text{safe}}$}
                              node[lab, below] {$v_{\text{scale}}$} (ctrl);
        \draw[arr] (ctrl) -- node[lab, above] {$v,\ \omega$} (veh);

        \draw[arr] (veh) -- node[lab, left, align=right] {noisy pose,\\speed} (est);
        \draw[arr] (est.west) -- ++(-1.5,0)
              node[lab, left, align=right] {$\hat{\mathbf{x}}$ to\\all blocks};
        \draw[arr] (det) -- (rep);
        \draw[arr] (rep) -- node[lab, right] {new path} (guid);
        \draw[arr] (det) -- node[lab, right] {detections} (avoid);
        \draw[arr] (lid) -- node[lab, right] {ranges} (avoid);
    \end{tikzpicture}
    \caption{System architecture. Shaded blocks are new in Part~B. The
    state estimate $\hat{\mathbf{x}}$ is an input to guidance, reactive
    avoidance, the controller and the re-planner; those four connections
    are drawn as one arrow for clarity.}
    \label{fig:architecture}
\end{figure*}

\begin{table*}[t]
    \centering
    \caption{Subsystem Inputs and Outputs}
    \label{tab:io}
    \footnotesize
    \begin{tabularx}{\textwidth}{@{}l X X l@{}}
        \toprule
        \textbf{Subsystem} & \textbf{Inputs} & \textbf{Outputs} & \textbf{Rate} \\
        \midrule
        Mission planner    & Wall map, start pose, five waypoints & Waypoint order, path & Once, before the mission \\
        State estimator    & Noisy pose and speed & State estimate $\hat{\mathbf{x}}$ & 100 Hz \\
        Guidance           & Path, waypoint order, $\hat{\mathbf{x}}$ & Heading command $\theta_{\text{cmd}}$, distance to waypoint $d_{\text{wp}}$, waypoint index & 100 Hz \\
        Re-planner         & $\hat{\mathbf{x}}$, detections, waypoint index & Replacement path & 100 Hz; re-plans at most once per second \\
        Reactive avoidance & $\theta_{\text{cmd}}$, $\hat{\mathbf{x}}$, detections, lidar ranges, $d_{\text{wp}}$ & Safe heading $\theta_{\text{safe}}$, speed factor $v_{\text{scale}}$ & 100 Hz \\
        Controller         & $\theta_{\text{safe}}$, $\hat{\theta}$, $d_{\text{wp}}$, $v_{\text{scale}}$ & Speed $v$ and turn-rate $\omega$ commands & 100 Hz \\
        \bottomrule
    \end{tabularx}
\end{table*}
